\section{Motivation}
\label{sec:motivation}

\subsection{Running Example: A Lost Wake-up Bug}
\label{subsec:running-example}

Figure~\ref{fig:running-rust} presents a typical synchronization defect in Rust involving a condition variable. In the \texttt{notifier} function, the notification (\texttt{notify\_one}) is invoked before the readiness flag is updated. If the \texttt{notifier} completes this sequence before the \texttt{worker} thread enters the \texttt{wait} state, the signal is lost, causing the \texttt{worker} to block indefinitely. Despite this critical liveness failure, the program adheres to Rust's ownership rules and compiles without warnings.

\begin{figure}[h]
\begin{lstlisting}[language=Rust, basicstyle=\small\ttfamily, frame=single]
fn worker(shared: Arc<(Mutex<bool>, Condvar)>) {
let (m, cv) = &*shared;
let mut g = m.lock().unwrap();
while !*g {
g = cv.wait(g).unwrap();
}
}

fn notifier(shared: Arc<(Mutex<bool>, Condvar)>) {
let (m, cv) = &*shared;
let mut g = m.lock().unwrap();
cv.notify_one();   // BUG: notify before flag update
*g = true;         // too late
}
\end{lstlisting}
\caption{A condition-variable signal-loss bug in Rust.}
\label{fig:running-rust}
\end{figure}

\subsection{The Complexity of Conventional Static Analysis}
\label{subsec:static-barriers}

Detecting this bug through traditional static analysis requires the simultaneous resolution of three interdependent challenges: identity of shared resources, implicit protection invariants, and non-deterministic interleaving. First, establishing the \textbf{identity of shared resources} requires proving that separate threads operate on the exact same synchronization primitives. In Rust, this necessitates tracing values through heap allocation, \texttt{Arc} cloning, and closure capture. Such pointer-aliasing problems frequently lead to a compromise between precision and scalability, often resulting in either \emph{false positives} due to over-approximation or \emph{false negatives} from under-approximation. 

Beyond the aliasing hurdle, implicit protection invariants such as the logical guarding of a boolean flag by a specific mutex are typically buried within implementation patterns rather than being explicitly declared. Traditional tools are often tightly coupled to specific language semantics, such as Rust’s \texttt{MutexGuard}, which restricts their generalizability across different memory models or languages. Even if these resource identities and invariants are successfully inferred, the analyzer must still perform an exhaustive interleaving analysis to expose the specific schedule where a notification is lost. Existing static analyzers frequently struggle at the initial stages of this pipeline. By shifting the focus from language-specific implementation artifacts to a language-agnostic abstraction of concurrent logic, we decouple the verification of synchronization patterns from the inherent complexities of pointer aliasing and source-level syntax.

\subsection{Bypassing Aliasing via \textsc{Cir} Abstraction}
\label{subsec:cir-abstraction}

The \textsc{Cir} in Figure~\ref{fig:running-cir} reframes the problem by shifting from implementation level inference to design level specification. Our architecture leverages an LLM to synthesize this alias-free representation from the program intent.

\begin{figure}[h]
\begin{lstlisting}[basicstyle=\small\ttfamily, frame=single]
resources:
m0:    { kind: Mutex, mode: Sync }
cv0:   { kind: Condvar, paired_with: m0 }
ready: { kind: Var, type: Bool, init: false }
protection:
ready: [m0]
functions:
worker:
body:
- { sid: w1, op: {lock: m0}, transfer: {next: w2} }
- { sid: w2, op: {read: ready}, transfer: {branch: {...}} }
- { sid: w3, op: {wait: {cv: cv0, mutex: m0}}, ... }
notifier:
body:
- { sid: n1, op: {lock: m0}, ... }
- { sid: n2, op: {notify_one: cv0}, transfer: {next: n3} }
- { sid: n3, op: {write: {var: ready, val: true}}, ... }
\end{lstlisting}
\caption{\textsc{Cir} for the running example, featuring global naming and explicit protection.}
\label{fig:running-cir}
\end{figure}

In \textsc{Cir}, shared identities are established by construction through global naming (\texttt{m0}, \texttt{cv0}). The protection relationship is elevated to a first-class declaration. By stripping away pointer indirection and heap-tracking, \textsc{Cir} transforms an intractable alias-analysis problem into a direct verification task on a clean, concurrent backbone.

\subsection{Synergistic Discovery and Repair}
\label{subsec:closing-loop}

The \textsc{Cir} is translated into a \textsc{Cvn}, which explores all interleavings. The formal engine identifies that if the transition at \texttt{n2} fires while the worker has not reached the wait place at \texttt{w3}, a signal loss occurs.

Crucially, the resulting counterexample is structured and anchored to specific statement IDs (\texttt{sid}). This allows the LLM to perform a targeted repair (e.g., swapping \texttt{n2} and \texttt{n3}) rather than a blind rewrite. This synergy exploits a fundamental division of labor: the LLM provides the semantic intuition to resolve resource identities, while the formal engine provides the exhaustive coverage required to expose edge-case interleavings.


